\pagenumbering{roman}
\setcounter{page}{1}

\selecthungarian

%----------------------------------------------------------------------------
% Abstract in Hungarian
%----------------------------------------------------------------------------
\chapter*{Kivonat}\addcontentsline{toc}{chapter}{Kivonat}

A mai világban egyre jobban terjed el, hogy az interneten elérhető szolgáltatások
 szerver alapú monolítikus architektúra helyett felhő alapú 
 infrastruktúrában kerülnek üzembehelyezésre.
 A Cloud Native Computing Foundation (CNCF) hivatalos definíciója \cite{CNCF} szerint
 a felhőnatív technológiák lehetővé teszik a szervezetek számára, hogy 
 skálázható alkalmazásokat  építsenek és futtassanak modern, dinamikus
 környezetekben, mint például a nyilvános, privát és hibrid felhők.
 A felhőnatív rendszerek során használt legfontosabb technológiák a konténerek, 
 a mikroszolgáltatások és az API (Application Programming Interface) -ok.
 Ezeknél a rendszereknél kiemelten fontos a rendelkezésre állás 
 és a rövid válaszidő mellett a megfelelő skálázódás biztosítása is.

 A jelen dokumentum a AI (Mesterséges Intelligencia) alapú ágensek használatát vizsgálja 
 felhős környezetekből kinyerhető metrikák alapján a skálázódás optimalizálására.
 A dokumentum bemutatja egy ilyen optimalizáló rendszer tervezését és megvalósítását, illetve kiértékelését
 különböző hatékonytalan forgatókönyvek alapján.







\vfill
\selectenglish


%----------------------------------------------------------------------------
% Abstract in English
%----------------------------------------------------------------------------
\chapter*{Abstract}\addcontentsline{toc}{chapter}{Abstract}

In today's world, it is more and more becoming a convention for web based services, that instead
 of server based monolithic architechture, they get deployed in a cloud based infrastructure.
 The Cloud Native Computing Foundation's official definition states that cloud native technologies give the
 ability to organizations to develop and deploy scalable applications in modern, dynamic envrionments, for example
 public, private and hybrid clouds.
 These cloud native systems typically utilize containers, microservices and APIs (Application Programming Interfaces).
 With these systems, ensuring proper scalability is paramount, alongside with high availability and low latency.

 This document investigates the use of AI (Artificial Intelligence) agents 
 for scaling optimization based on metrics that describe the cloud environment.
 The document showcases the design and implementation of such an optimization system 
 and evaluates its performance under various inefficient scenarios.
    



\vfill
\selectthesislanguage

\newcounter{romanPage}
\setcounter{romanPage}{\value{page}}
\stepcounter{romanPage}